\chapter{Configuration Model and Validation}
\label{chap:configuration-validation}

This chapter describes the configuration layer of the framework, that is, the
representation used to describe a network laboratory before any Kathar\'a file
is generated. The framework adopts YAML as a declarative source format
\cite{yaml122} and separates the description of the desired topology from the
commands required to realise it. A single configuration can therefore contain the laboratory
identity, nodes, interfaces, IP addressing, static routes, optional startup
commands, and dynamic-routing parameters. The same source is then reused both
for validation and, after validation has succeeded, for generation.

The implementation discussed here is contained mainly in
\texttt{src\_gen\_lab/Lexer.py} and \texttt{src\_gen\_lab/Parser.py}. The
historical filenames are retained in the repository, although their practical
responsibilities are closer to configuration loading and validation than to a
traditional compiler front end. The YAML document is loaded as nested Python
dictionaries and lists; the framework does not construct a separate domain
class for every node, interface, or routing protocol. Consequently, the
accepted model is defined by the fields interpreted by the generator and by
the structural and semantic checks implemented by the parser.

The chapter stops at the validated configuration. The creation of
\texttt{lab.conf}, startup files, FRRouting files, optional monitoring
components, and the laboratory lifecycle are discussed in the following
chapter.

\section{Configuration Model and Supported Operations}
\label{sec:configuration-model-purpose}

A configuration can be selected either by passing an explicit YAML path or by
passing only its short name. The function \texttt{resolve\_config\_path()} first
checks whether the supplied argument already identifies an existing file. If
it does not, the function searches the configured directory, which defaults to
\texttt{configs/}, and tries the extensions \texttt{.yml} and \texttt{.yaml}.
This permits both forms shown in Listing~\ref{lst:config-resolution}.

\begin{lstlisting}[language=bash,caption={Equivalent ways of selecting a configuration.},label={lst:config-resolution}]
python3 generate_lab.py configs/client-router-server.yml
python3 generate_lab.py client-router-server
\end{lstlisting}

Once the path has been resolved, \texttt{load\_yaml()} opens the file using
UTF-8 encoding and calls \texttt{yaml.safe\_load()}, following the safe-loading
mechanism documented by PyYAML \cite{pyyaml_docs}. Safe loading is important
for a configuration file because it avoids the general object-construction
behaviour associated with unsafe YAML loading. The loader also checks that the top-level value is a
Python dictionary; an empty value, scalar, or list is rejected before the
network-specific validation begins. If PyYAML is unavailable, the loader
reports the missing dependency explicitly.

Every laboratory configuration must contain at least \texttt{lab\_name} and
\texttt{nodes}. The remaining fields are optional and are included only when
the corresponding behaviour is needed. A compact top-level example is shown
in Listing~\ref{lst:top-level-yaml}.

\begin{lstlisting}[caption={General top-level structure of a laboratory configuration.},label={lst:top-level-yaml}]
lab_name: example_lab
description: Example routed laboratory

metadata:
  description: "Example generated with the framework"
  version: 2.0
  author: "Luca-Bortoluzzi"
  email: "luca.bortoluzzi921@edu.unito.it"
  web: "http://www.kathara.org/"

nodes:
  # node declarations
\end{lstlisting}

The \texttt{lab\_name} value identifies the generated laboratory and becomes
the directory name under the selected output directory. The optional
\texttt{description} and \texttt{metadata} mapping do not affect forwarding or
addressing. They provide descriptive values later written to
\texttt{lab.conf}; when a metadata field is omitted, the generator uses a
fallback. The \texttt{nodes} field is a non-empty mapping in which each key is
a node name and each value is the corresponding node declaration. Node names
and collision-domain names must satisfy the same identifier rule:

\begin{lstlisting}[caption={Identifier pattern accepted for nodes and collision domains.},label={lst:name-pattern}]
^[A-Za-z0-9][A-Za-z0-9_.-]*$
\end{lstlisting}

Therefore, letters, digits, dots, underscores, and hyphens are accepted, while
spaces and identifiers beginning with an invalid character are rejected. The
node name is not only descriptive: it is reused in generated filenames,
\texttt{lab.conf}, and, when node resources are imported, in the matching node
directory.

Listing~\ref{lst:annotated-node-reference} shows a representative router and
annotates the principal fields. The example deliberately combines several
features in one declaration; a normal node generally uses only the subset that
is relevant to its role.

\begin{lstlisting}[caption={Annotated node configuration reference.},label={lst:annotated-node-reference}]
nodes:
  r1:                              # node name
    type: router                   # host | router | switch | firewall
    image: kathara/frr             # optional image override
    forwarding: true               # forwarding-consistency declaration

    interfaces:                    # ordered list: first entry -> eth0
      - network: lan_a             # collision-domain name
        ip: 10.10.1.1/24           # address in CIDR notation
        ospf_area: 10              # optional OSPF area

      - network: r1_r2             # second entry -> eth1
        ip: 10.0.12.1/30
        rip: true                  # optional RIP marker

    default_gateway: 10.10.1.254   # optional default-route next hop

    routes:                        # optional static routes
      - to: 192.168.20.0/24
        via: 10.0.12.2

    commands:                      # optional startup commands
      - sysctl -w net.ipv4.ip_forward=1
      - echo "r1 initialised"

    frr:                           # optional FRRouting configuration
      enabled: true
      protocols: [ospf, bgp]
      router_id: 1.1.1.1

      ospf:
        areas:
          - id: 10
            type: stub
            no_summary: false
        redistribute: [bgp]

      bgp:
        asn: 65001
        networks:
          - 10.10.1.0/24
        neighbors:
          - ip: 10.0.12.2
            remote_as: 65002
        redistribute: [ospf]
        ebgp_requires_policy: false
        network_import_check: false
\end{lstlisting}

Table~\ref{tab:supported-configuration-operations} is the operational reference
for node configuration. Each row corresponds to one operation that the user
can express in YAML and states what the framework does with it. Related nested
fields are grouped in the same row so that the complete model remains readable
without turning each property into a separate subsection.

{\small
\begin{longtable}{>{\raggedright\arraybackslash}p{0.20\textwidth}>{\raggedright\arraybackslash}p{0.24\textwidth}>{\raggedright\arraybackslash}p{0.44\textwidth}}
\caption{Supported configuration operations and their effect.}
\label{tab:supported-configuration-operations}\\
\toprule
\textbf{Operation} & \textbf{YAML field or syntax} & \textbf{Effect in the framework} \\
\midrule
\endfirsthead
\toprule
\textbf{Operation} & \textbf{YAML field or syntax} & \textbf{Effect in the framework} \\
\midrule
\endhead
Identify the laboratory & \texttt{lab\_name} & Mandatory top-level value. It identifies the laboratory and is later used as the generated directory name. \\
\\
Add descriptive metadata & \texttt{description}, \path{metadata.description}, \path{metadata.version}, \path{metadata.author}, \path{metadata.email}, \path{metadata.web} & Adds human-readable information that is later transferred to \texttt{lab.conf}. These values do not change the topology or routing behaviour. \\
\\
Create and name a node & \path{nodes.<name>} & Creates the node identifier. The name is reused in Kathar\'a declarations, startup filenames, and optional node resource directories. \\
Select the node role and image & \texttt{type}, \texttt{image} & \texttt{type} accepts \texttt{host}, \texttt{router}, \texttt{switch}, or \texttt{firewall}; omitted types are treated as hosts. \texttt{image} overrides automatic Docker-image selection. Ordinary hosts default to \texttt{kathara/base}; router, switch, and firewall nodes are treated as FRR-capable and default to \texttt{kathara/frr}. \\
\\
Declare forwarding intent & \texttt{forwarding} & Provides a role-consistency declaration. A router explicitly configured with \texttt{forwarding: false} is rejected by semantic validation. The field does not itself generate a kernel \texttt{sysctl} command. \\
\\
Create interfaces & \texttt{interfaces} & Mandatory non-empty ordered list. Entry zero becomes \texttt{eth0}, entry one becomes \texttt{eth1}, and so on. This ordering is later reused both in \texttt{lab.conf} and startup generation. \\
\\
Connect an interface to a network & \path{interfaces[].network} & Mandatory for every interface. The value is the Kathar\'a collision-domain name. Interfaces on different nodes using the same value belong to the same Layer~2 segment. \\
\\
Assign an interface address & \path{interfaces[].ip} & Optional IPv4 or IPv6 interface in CIDR notation. During generation it becomes an \texttt{ip address add ... dev ethN} command. An interface can remain unnumbered when only Layer~2 attachment is required. \\
\\
Mark interfaces for dynamic routing & \path{interfaces[].ospf_area}, \path{interfaces[].rip} & \texttt{ospf\_area} associates the interface network with an OSPF area and requires an address. \texttt{rip: true} marks the interface network for RIP; if no interface is marked, RIP generation can fall back to all addressed interfaces. \\
\\
Configure a default route & \texttt{default\_gateway} & Optional next-hop address. During generation it becomes \texttt{ip route add default via ...}. Semantic validation requires the gateway to belong to one of the node's directly connected subnets. \\
\\
Configure static routes & \texttt{routes}, \path{routes[].to}, \path{routes[].via} & Adds explicit routes to the startup file. Each mapping states a destination prefix in \texttt{to} and a next hop in \texttt{via}. The current parser does not perform the same depth of semantic analysis on arbitrary static routes as it does on interfaces and gateways. \\
\\
Append custom startup commands & \texttt{commands} & Appends the supplied strings to the startup file in declaration order. This is the generic escape mechanism for kernel settings, service startup, file manipulation, or other operations not represented by dedicated YAML fields. Arbitrary command semantics are not generally validated. \\
\\
Activate FRRouting and select protocols & \texttt{frr}, \texttt{frr.enabled}, \texttt{frr.protocol}, \texttt{frr.protocols}, \texttt{frr.daemons} & Declares FRRouting behaviour. A single protocol can be selected with \texttt{protocol}; multiple protocols can be listed or configured through \texttt{protocols}; daemon flags can also imply OSPF, RIP, or BGP. The supported protocols are OSPF, RIP, and BGP. \\
\\
Assign routing identifiers & \texttt{frr.router\_id}, \path{frr.ospf.router_id}, \path{frr.bgp.router_id} & Sets a common router ID or a protocol-specific override. Router IDs are parsed as IPv4 addresses and must be unique across the effective declarations checked by the validator. \\
\\
Configure OSPF & \texttt{frr.area}, \path{frr.ospf.networks}, \path{frr.ospf.areas}, \path{frr.ospf.areas[].id}, \path{frr.ospf.areas[].type}, \path{frr.ospf.areas[].no_summary}, \path{frr.ospf.areas[].totally_stub}, \path{frr.ospf.redistribute} & Supports the legacy common-area form, interface-level areas, explicit networks, stub or NSSA area declarations, optional no-summary behaviour, and redistribution. If interface-level \texttt{ospf\_area} markers exist, they determine the generated network-area pairs; otherwise the generator falls back to the common area and FRR network list. \\
\\
Configure RIP & \path{frr.rip.version}, \path{frr.rip.networks}, \path{frr.rip.auto_summary}, \path{frr.rip.redistribute} & Selects the RIP version, network statements, automatic summarisation behaviour, and redistribution. When no explicit network list is supplied, the generator uses RIP-marked interfaces or, as a fallback, all addressed interfaces. \\
\\
Configure BGP & \path{frr.bgp.asn}, \path{frr.bgp.router_id}, \path{frr.bgp.networks}, \path{frr.bgp.neighbors}, \path{neighbors[].ip}, \path{neighbors[].remote_as}, \path{frr.bgp.redistribute}, \path{frr.bgp.ebgp_requires_policy}, \path{frr.bgp.network_import_check} & Defines the local AS, router ID, announced prefixes, peers and their remote AS values, redistribution, and the FRRouting eBGP-policy and network-import switches. The ASN must be in the range 1--4294967295. Peer addresses and remote AS values are checked against nodes declared in the same laboratory. \\
\\
Provide a complete FRR configuration & \texttt{frr.config} & Replaces automatic protocol-block generation with user-supplied FRRouting text written directly to \texttt{frr.conf}. It is intended for advanced configurations that exceed the declarative model. \\
\bottomrule
\end{longtable}
}

The node types share this common syntax; the framework does not define a
separate YAML language for hosts, routers, switches, and firewalls. Their
practical differences derive mainly from validation expectations and automatic
image selection. A host normally contains one or more interfaces, an optional
default gateway, and possibly startup commands. A router normally contains at
least two addressed interfaces and either static routes or an FRRouting
configuration. In the current implementation, \texttt{switch} is only a role
label used by validation and image selection: it does not create a Linux
bridge, Open vSwitch instance, or any other Layer~2 forwarding function.
Kathar\'a's \texttt{device[n]} syntax only attaches \texttt{ethn} to the named
collision domain \cite{kathara_labconf_manpage}; consequently, two interfaces
on a node do not by themselves join two collision domains. A functioning
Layer~2 switch requires explicit bridge or Open vSwitch configuration, for
example through imported resources and startup commands or a suitable image
such as \texttt{kathara/openvswitch} \cite{kathara_images_repo}. A firewall
uses the same interface model and currently relies on generic commands or
imported resources for policy-specific behaviour because the YAML model does
not define a dedicated firewall-policy language.

The \texttt{image} field allows these defaults to be overridden. If it is
absent, ordinary hosts use \texttt{kathara/base}; nodes whose type is
\texttt{router}, \texttt{switch}, or \texttt{firewall} are considered
FRR-capable and use \texttt{kathara/frr}. A host can also become FRR-capable by
providing an FRR declaration. One implementation detail deserves attention:
for router, switch, and firewall types the device type itself is sufficient to
select the FRR execution path. Therefore, \texttt{frr.enabled} should not be
interpreted as a universal switch that disables FRR for every node type; its
behaviour is most meaningful for a host whose FRR capability comes from the
\texttt{frr} section. This is one of the implementation details that a future
revision could make more uniform.

Interfaces are the core of the topology model. Each node must contain a
non-empty \texttt{interfaces} list and every interface must have a
\texttt{network} field. The network name represents a Kathar\'a collision
domain, while the optional \texttt{ip} field assigns an IPv4 or IPv6 address in
CIDR notation. List order is operational: the first interface becomes
\texttt{eth0}, the second \texttt{eth1}, and so forth. The same order is later
used when the generator creates \texttt{lab.conf} and node startup commands.
Listing~\ref{lst:host-yaml} shows the minimal pattern used by a routed host.

\begin{lstlisting}[caption={Host with one addressed interface and a default gateway.},label={lst:host-yaml}]
pc_a:
  type: host
  interfaces:
    - network: lan_a
      ip: 10.10.1.10/24
  default_gateway: 10.10.1.1
\end{lstlisting}

An interface may omit its address, which is useful for an unnumbered or
Layer~2-only attachment. Protocol-specific fields can impose stronger
requirements: for example, an interface carrying \texttt{ospf\_area} must also
have an IP address because the generator derives the advertised prefix from
that address. The optional default gateway is validated as an IP address and
must fall inside at least one locally connected subnet, which prevents a
syntactically valid but Layer~2-unreachable gateway from reaching generation.
Additional routes are declared through \texttt{routes}; every route contains a
\texttt{to} destination and a \texttt{via} next hop. The generator translates
these mappings directly into \texttt{ip route add} commands. The current
validator intentionally devotes more semantic analysis to interfaces,
subnets, gateways, and dynamic-routing relationships than to arbitrary static
route policies, so malformed or logically inappropriate custom route content
is less extensively checked.

The \texttt{commands} list was introduced as a pragmatic escape mechanism for
behaviour outside the formal model. During development, not every node-specific
operation justified a new dedicated YAML field: application startup, temporary
kernel settings, and scenario-specific file operations could otherwise have
forced the configuration language to grow around individual experiments. Each
command string is therefore copied to the node startup file in declaration
order. This choice kept the core model small and allowed specialised scenarios
to be implemented without blocking on a new parser feature.

The trade-off is that command text is opaque to the general parser, so the
framework cannot prove its correctness. It performs only limited consistency
checks, such as detecting an explicit \path{ip_forward=0} setting on a node
declared as a router. Similarly, \texttt{forwarding: true} is a declaration of
intent and validation consistency; it does not by itself emit a forwarding
command. When explicit kernel configuration is required, it must be expressed
through \texttt{commands} or provided by the selected container environment.
Chapter~\ref{chap:robustness-security-scalability} later shows why this useful
development shortcut also creates a trust and robustness boundary.

Dynamic routing was added as the laboratories evolved beyond static
host--router scenarios. Supporting OSPF, RIP, and BGP required the
configuration language to represent protocol-specific information that was not
needed by the first prototypes. The model was therefore extended incrementally
instead of being replaced by a new syntax. For compatibility with
configurations developed at different stages of the project, the generator
accepts several declaration forms. A single protocol can be written as
\texttt{frr.protocol}; multiple protocols can be placed in a list under
\texttt{frr.protocols}; \texttt{frr.protocols} can also be a mapping containing
protocol-specific options; dedicated \texttt{frr.ospf}, \texttt{frr.rip}, and
\texttt{frr.bgp} mappings are recognised; and daemon flags under
\texttt{frr.daemons} can imply the corresponding protocol. These forms are
illustrated in Listing~\ref{lst:frr-declaration-forms}.

\begin{lstlisting}[caption={Compatible FRRouting protocol declaration forms.},label={lst:frr-declaration-forms}]
# One protocol
frr:
  enabled: true
  protocol: ospf

# Multiple protocols
frr:
  protocols: [ospf, bgp]

# Structured multiprotocol form
frr:
  protocols:
    ospf:
      enabled: true
      redistribute: [bgp]
    bgp:
      enabled: true
      asn: 65001

# Direct daemon selection retained for compatibility
frr:
  daemons:
    ospfd: true
    ripd: false
    bgpd: true
\end{lstlisting}

Retaining these forms was a deliberate compatibility decision. Requiring an
immediate migration to one final representation would have invalidated earlier
test topologies precisely when they were useful for regression testing. The
cost is a broader parser and some representational redundancy. For a mature
configuration language, a single versioned schema would be preferable; the
strict-schema direction discussed in
Chapter~\ref{chap:robustness-security-scalability} is intended to reduce this
ambiguity without losing the lessons provided by the earlier files.

OSPF can use a common legacy area declared through \texttt{frr.area}, or each
interface can carry its own \texttt{ospf\_area}. The latter form is used by the
multi-area examples because it keeps area membership next to the interface
whose network will be advertised. When one or more interface area markers are
present, the generator derives network-area pairs from them. If no such marker
exists, the generator falls back to the FRR network list and the common area,
which defaults to area~0. The structured OSPF mapping can also contain
explicit networks, area declarations, stub and NSSA area types, optional
\texttt{no\_summary} or \texttt{totally\_stub} flags, and redistribution.
The underlying OSPFv2 and NSSA concepts are defined in RFC~2328 and RFC~3101
\cite{rfc2328,rfc3101}.
Listing~\ref{lst:ospf-yaml} shows a concise multi-area example.

\begin{lstlisting}[caption={Interface-level OSPF area configuration.},label={lst:ospf-yaml}]
r1:
  type: router
  interfaces:
    - network: lan_a
      ip: 10.10.1.1/24
      ospf_area: 10
    - network: backbone
      ip: 10.0.12.1/30
      ospf_area: 0
  frr:
    enabled: true
    protocol: ospf
    router_id: 1.1.1.1
\end{lstlisting}

RIP follows the same general model and follows the RIPv2 concepts defined in
RFC~2453 \cite{rfc2453}. It can be selected through the common
protocol field or a dedicated \texttt{frr.rip} mapping. The configuration can
choose the protocol version, control automatic summarisation, provide an
explicit network list, and request redistribution. Interface-level
\texttt{rip: true} markers provide a concise way to select the participating
networks. If no marked interfaces or explicit networks are provided, the
generator can derive RIP networks from all addressed interfaces.

BGP requires additional relationship information consistent with the BGP-4
model defined in RFC~4271 \cite{rfc4271}. The local autonomous-system
number can be written at the common FRR level or inside the BGP mapping; the
structured value takes precedence. BGP can also define a router ID, announced
networks, neighbours, redistribution, and two switches used by the generated
FRRouting configuration: \texttt{ebgp\_requires\_policy} and
\texttt{network\_import\_check}. The former is directly related to the
default external-BGP policy behaviour specified by RFC~8212
\cite{rfc8212}. Each neighbour contains an \texttt{ip} and a
\texttt{remote\_as}. Listing~\ref{lst:bgp-yaml} shows the compact form used in
the multi-AS configuration.

\begin{lstlisting}[caption={BGP node declaration.},label={lst:bgp-yaml}]
r1:
  type: router
  interfaces:
    - network: lan_a
      ip: 10.10.1.1/24
    - network: r1_r2
      ip: 10.0.12.1/30
  frr:
    enabled: true
    protocol: bgp
    asn: 65001
    router_id: 1.1.1.1
    networks:
      - 10.10.1.0/24
    neighbors:
      - ip: 10.0.12.2
        remote_as: 65002
\end{lstlisting}

A common or protocol-specific router ID is interpreted as an IPv4 identifier.
The common value under \texttt{frr.router\_id} is used unless the OSPF or BGP
section provides its own override. For configurations that cannot be expressed
conveniently through these fields, \texttt{frr.config} acts as a final escape
hatch: its complete multiline value replaces automatically generated protocol
blocks and is written directly to \texttt{frr.conf}. This preserves
extensibility without requiring the YAML model to represent every FRRouting
command.

Not every key that may appear in a repository YAML file is part of the core
generation language. YAML loading preserves unknown mappings because the
configuration remains an ordinary Python dictionary, but only fields
explicitly interpreted by the generator have a generation effect. For
example, a scenario can contain descriptive or auxiliary sections, and some
example files contain additional information used by separate tools. Such data
is not automatically converted into Kathar\'a configuration simply because it
is present in the YAML document. This distinction is important when reading
the repository: the accepted YAML document can be broader than the core node
model documented in Table~\ref{tab:supported-configuration-operations}.

\section{Validation Process and Semantic Rules}
\label{sec:structural-validation}

Validation is divided into two stages because malformed structure and
network-level inconsistency require different handling. Structural validation
is implemented by \texttt{validate\_config()} and is intentionally fail-fast.
It verifies the assumptions that must hold before more detailed analysis can
run: \texttt{lab\_name} must exist; \texttt{nodes} must exist and be a
non-empty dictionary; every node declaration must itself be a dictionary;
every node must contain a non-empty interface list; every interface entry must
be a dictionary; and every interface must contain \texttt{network}. A
violation raises \texttt{ValueError} immediately. This prevents a missing
basic field from producing secondary implementation errors during semantic
analysis. For example, collision-domain grouping cannot produce a meaningful
diagnostic if an interface is not a mapping or has no network name.

Semantic validation is implemented by \texttt{validate\_semantics()} and uses a
different policy: rather than stopping after the first network-level problem,
it collects all detectable issues in one \texttt{SemanticValidationReport}.
Each issue is represented by an immutable Python dataclass
\cite{python_dataclasses}, shown in Listing~\ref{lst:validation-issue}.

\begin{lstlisting}[language=Python,caption={Representation of a semantic validation issue.},label={lst:validation-issue}]
@dataclass(frozen=True)
class ValidationIssue:
    severity: str
    code: str
    message: str
\end{lstlisting}

The severity is \texttt{ERROR} or \texttt{WARN}. The code provides a stable
machine-oriented category and the message gives the human-readable context.
The report stores the node, interface, and collision-domain counts, records the
number of OSPF and BGP nodes, tracks whether resource directories were checked,
and holds the complete issue list. Computed properties expose the filtered
errors and warnings and a \texttt{valid} flag. A report is valid when it
contains no error; warnings do not block generation. This distinction allows
unusual but potentially intentional topologies to proceed while still drawing
attention to suspicious choices.

The first semantic rules concern identifiers and node roles. Node names and
collision-domain names are checked against the regular expression introduced
in Section~\ref{sec:configuration-model-purpose}. The supported node types are
\texttt{host}, \texttt{router}, \texttt{switch}, and \texttt{firewall}; any
other value is an error. Routers receive two additional consistency checks. If
\texttt{forwarding} is explicitly false, or if the custom command list
contains a setting equivalent to \path{ip_forward=0}, the declaration is
rejected because the node role conflicts with its stated forwarding
behaviour. A router with only one interface is not structurally impossible,
but it cannot forward between two network segments, so the validator reports a
warning rather than an error.

Interface addresses are parsed with Python's \texttt{ipaddress.ip\_interface()}
function \cite{python_ipaddress}. Both IPv4 and IPv6 interface notation are
accepted by this generic parser. Malformed values generate an error. For IPv4 networks with a prefix
shorter than \texttt{/31}, assigning the network address or broadcast address
to an interface is also rejected. Every valid host address is indexed across
the complete laboratory, allowing duplicate addresses to be reported with all
owning node-interface pairs. The implemented rule is global: two isolated
segments are not allowed to reuse the same address, even if they would never
communicate. This reflects the framework's assumption of one coherent address
space per laboratory and simplifies subsequent target and topology handling.

Collision domains are then checked as Layer~2 segments. The validator groups
all endpoints sharing the same \texttt{network} name. If a domain has only one
endpoint, the condition is reported as a warning because it may be deliberate,
but it often indicates a misspelled network name or a missing peer. Addressed
interfaces belonging to the same collision domain are expected to use one
coherent subnet; if more than one distinct subnet appears on the same domain,
the configuration is rejected. The validator also compares the subnets used by
different collision domains and rejects overlapping prefixes. This avoids a
topology in which separate Layer~2 segments are represented by identical or
ambiguous Layer~3 address spaces.

The default gateway is checked after interface parsing. The value must be a
valid IP address and must belong to at least one local interface subnet of the
same node. A gateway that is syntactically valid but outside all locally
connected networks is therefore rejected before generation. This rule is
particularly useful for host declarations, where a wrong gateway can otherwise
produce a laboratory that starts correctly while being unable to reach any
remote subnet.

Some additional role-oriented checks exist in the current parser because the
repository contains auxiliary scenarios built on top of the same node model.
These checks require an address for certain reserved role naming conventions.
They are implementation-specific conventions rather than a general property of
Kathar\'a or of the four node types, and they do not change the generic
configuration operations documented in this chapter.

FRRouting validation begins by normalising the different compatible protocol
declaration forms. The helper that collects declared protocols accepts
\texttt{protocol}, a list or mapping under \texttt{protocols}, the dedicated
\texttt{ospf}, \texttt{rip}, and \texttt{bgp} mappings, and daemon flags. Any
protocol outside OSPF, RIP, and BGP is reported as unsupported. The parser also
requires the \texttt{frr} section, when present and active, to be a dictionary.
Effective router IDs are collected from the common field and from
protocol-specific overrides. Each router ID must parse as an IPv4 address, and
duplicate effective IDs across nodes are rejected.

For OSPF, an interface marked with \texttt{ospf\_area} must contain an IP
address. Explicit OSPF network values are parsed as prefixes and must overlap at
least one network directly connected to the node. This detects stale or copied
network statements that do not correspond to the actual interface addressing.
The check is deliberately local: it verifies that the node's declaration is
coherent, but it does not attempt to prove global area connectivity or protocol
convergence.

For BGP, the local autonomous-system number must be an integer from 1 to
4294967295. Neighbours must be expressed as a list of mappings. Each neighbour
IP must be valid, must belong to an interface of a declared node, and cannot
refer back to the local node itself. The neighbour's \texttt{remote\_as} must
also be valid and must match the ASN configured on the peer node. If the peer
has no valid BGP ASN, the relationship is rejected. After these mandatory
checks, the parser also examines reciprocity: when node A declares an address
of node B as a neighbour but node B does not declare any address of node A, a
warning is emitted. Reciprocity is treated as a warning rather than an error
because an incomplete adjacency may be intentional while a topology is being
constructed incrementally.

The validator also performs reference-integrity checks for auxiliary
configuration data stored in the same YAML document. A visible generic example
is the mapping used for connection tests, named \texttt{connection\_tests} in
the YAML files. A declared client must refer to an existing node, a target node
must exist and have an address, and a target IP must be valid and assigned to a
declared interface. These checks do not execute a test and do not make the
auxiliary subsystem part of laboratory generation; they merely ensure that
references embedded in the source configuration remain consistent with the
topology. Other auxiliary sections can be checked through the same internal
reference mechanism without altering the node model.

When node-directory import validation is requested, the parser examines
visible directories beside the selected YAML file. Directories such as
\path{__pycache__} and hidden entries are excluded. A visible directory whose
name does not correspond to a declared node is reported as a warning and will
be skipped by the generator. This policy is useful because the shared
\texttt{configs/} directory can contain resources for several laboratories at
the same time; an unrelated resource directory should be reported but should
not automatically invalidate an otherwise correct topology.

Table~\ref{tab:error-warning-policy} summarises the principal rules. It
combines the checks described above without introducing a separate subsection
for every category.

{\small
\begin{longtable}{>{\raggedright\arraybackslash}p{0.37\textwidth}>{\raggedright\arraybackslash}p{0.12\textwidth}>{\raggedright\arraybackslash}p{0.39\textwidth}}
\caption{Representative validation rules and their consequences.}
\label{tab:error-warning-policy}\\
\toprule
\textbf{Condition} & \textbf{Severity} & \textbf{Consequence} \\
\midrule
\endfirsthead
\toprule
\textbf{Condition} & \textbf{Severity} & \textbf{Consequence} \\
\midrule
\endhead
Missing \texttt{lab\_name}, missing or empty \texttt{nodes}, malformed node or interface structure, or missing interface \texttt{network} & Error & Structural validation stops immediately because later analysis cannot operate safely. \\
Invalid node or collision-domain identifier & Error & Generation is blocked. \\
Unsupported node type & Error & Generation is blocked. \\
Malformed interface address & Error & Generation is blocked. \\
IPv4 network or broadcast address assigned to an ordinary host interface prefix shorter than /31 & Error & Generation is blocked. \\
Duplicate interface address anywhere in the laboratory & Error & Generation is blocked. \\
Collision domain with only one endpoint & Warning & Generation may continue; the domain may be deliberately isolated or may contain a naming mistake. \\
Different subnets used on the same collision domain & Error & Generation is blocked because one Layer~2 segment has inconsistent Layer~3 addressing. \\
Overlapping subnets used on different collision domains & Error & Generation is blocked because the address model is ambiguous. \\
Invalid default gateway or gateway outside every local subnet & Error & Generation is blocked. \\
Router with one interface & Warning & Generation may continue, but the declared router cannot forward between distinct segments. \\
Router forwarding explicitly disabled & Error & Generation is blocked because the declaration contradicts the router role. \\
Invalid, duplicated, or unsupported FRRouting declaration & Error & Generation is blocked when the affected routing configuration cannot be considered coherent. \\
OSPF area marker without an IP address or explicit OSPF network disconnected from every local interface & Error & Generation is blocked. \\
Invalid BGP ASN, invalid peer address, self-neighbour, missing peer node, or remote-AS mismatch & Error & Generation is blocked. \\
BGP adjacency declared in only one direction & Warning & Generation may continue while the incomplete or intentionally asymmetric declaration is reviewed. \\
Auxiliary target or client reference to a missing node or undeclared address & Error & The source configuration is inconsistent and generation is blocked. \\
Resource directory without a matching node when import checking is enabled & Warning & The directory is reported and skipped; the topology can still be valid. \\
\bottomrule
\end{longtable}
}

A valid report can therefore contain warnings. The \texttt{valid} property is
based only on the absence of errors. The report's helper methods can append
errors and warnings and determine whether specific error-code prefixes are
present, which allows the textual summary to suppress an \texttt{[OK]} message
when a related error already exists. A successful run can produce output of
the form shown in Listing~\ref{lst:semantic-report-output}.

\begin{lstlisting}[caption={Example of a successful semantic validation report.},label={lst:semantic-report-output}]
[OK] 9 nodes validated
[OK] 8 collision domains validated
[OK] IP addresses are valid and unique
[OK] Subnets and default gateways are coherent
[OK] OSPF configuration is coherent on 5 nodes
[OK] Semantic validation passed with 0 warning(s)
\end{lstlisting}

This two-stage design is useful from a usability perspective and also reflects
an implementation problem observed during development. Structural errors must
be reported immediately because the document cannot be interpreted safely, but
stopping semantic analysis after its first error made larger topologies slow to
correct: fixing one address or routing inconsistency could simply reveal the
next independent problem on the following run. Semantic validation was
therefore designed to collect multiple issues in one report. A user can correct
duplicated addresses, an incorrect gateway, and a BGP remote-AS mismatch
together instead of discovering them through three consecutive generator
failures. In this sense, the report format is not only diagnostic output; it is
a direct optimisation of the edit--validate--correct development cycle.

\section{Validation Workflow and Limitations}
\label{sec:validation-workflow}

The command-line option \cliopt{validate-only} exposes the validation layer
without creating or modifying a laboratory. The workflow resolves the
configuration path, loads the YAML document, executes structural validation,
executes semantic validation, prints the resulting report, and terminates
before output-directory preparation. A typical invocation is shown in
Listing~\ref{lst:validate-only}.

\begin{lstlisting}[language=bash,caption={Validating a configuration without generating a laboratory.},label={lst:validate-only}]
python3 generate_lab.py multi-area --validate-only
\end{lstlisting}

This mode is useful while editing a topology because validation can be repeated
without changing generated files. It can also be incorporated into a
version-control or continuous-integration workflow. Normal generation performs
semantic validation again before writing output. The repetition is
intentional: the generation function should not rely on being called only
through one command-line path, so it protects itself by validating the
configuration it receives.

The repository contains several configurations of increasing complexity that
exercise the same core model. \texttt{client-router-server.yml} represents a
minimal routed topology and is useful for checking the fundamental mapping of
hosts, interfaces, collision domains, and default gateways. \texttt{Test1.yml}
contains several client LANs, a server-side network, and a mesh of OSPF routers;
it exercises a larger single-area addressing and routing configuration and
also contains auxiliary connection-test references that can be checked for
integrity. \texttt{multi-area.yml} associates OSPF areas directly with
interfaces and demonstrates that the area assignment can remain close to the
prefix it affects. \texttt{multi\_AS.yml} assigns different BGP autonomous
systems to routers and declares explicit peers and advertised LAN prefixes,
thereby exercising ASN, neighbour-address, remote-AS, reciprocity, and router-ID
checks. Finally, \texttt{Test\_enterprise\_multi\_as\_structured.yml} combines
several routing styles and uses the structured multiprotocol representation,
including OSPF area options, RIP markers, BGP sessions, and redistribution.
Together, these files act both as examples of the accepted language and as
regression inputs during framework development.

The current configuration model remains intentionally pragmatic. It is not
formalised through a separate YAML Schema or JSON Schema document; the Python
implementation is the authoritative definition of accepted structure and
semantics. This keeps the project compact, but external editors cannot derive
the complete language from a standard schema file. The configuration also
remains a hierarchy of dictionaries and lists rather than a typed internal
object model, so some assumptions are distributed across the loader, validator,
and generator.

Not all aspects of a node can be validated statically. Custom startup commands
are arbitrary shell strings, and a complete manual \texttt{frr.config} block is
opaque to most semantic checks. The framework can validate the surrounding
node and interface model, but the correctness of user-supplied command text is
ultimately the user's responsibility. Static route entries are similarly less
richly analysed than addressing and dynamic-routing relationships. A future
version could define stricter schemas for these structures while retaining a
controlled escape mechanism for advanced use cases.

Some current rules depend on project conventions rather than on a universal
network model. The role-oriented address checks mentioned above rely on node
naming conventions. The global duplicate-address rule also assumes one address
namespace for the complete laboratory; supporting overlapping addresses in
separate routing instances would require a richer representation of routing
contexts. The resource-directory check relies on exact correspondence between
a directory name and a declared node. These choices are reasonable for the
current framework because they simplify generation and diagnostics, but they
should be understood as implementation constraints rather than Kathar\'a
requirements.

The FRRouting model is intentionally broad enough for the topologies developed
in the project but does not attempt to expose the entire FRRouting command
language. OSPF, RIP, and BGP are supported, including multiprotocol routers and
redistribution, while advanced features can be supplied through
\texttt{frr.config}. As noted earlier, FRR capability is partly inferred from
node type and partly from the \texttt{frr} section. This produces a small
asymmetry around \texttt{frr.enabled} for nodes that are already typed as
routers, switches, or firewalls. Making capability selection explicit and
uniform would be a useful refinement.

Another important limit is that semantic validation proves configuration
coherence, not runtime correctness. A YAML file can have valid identifiers,
unique addresses, consistent subnets, reachable gateways, and coherent routing
relationships while the resulting network still fails to converge because of
a runtime condition, a user-supplied command, a container-specific issue, or a
more subtle protocol property not represented by the validator. Application
availability is even further outside the scope of this stage. These properties
must be checked after generation and execution. The value of the validation
layer is therefore preventive: it eliminates a large class of deterministic
configuration mistakes before containers are started, reducing the number of
runtime failures that must be diagnosed manually.

In summary, the YAML document is the authoritative description of a laboratory
before generation. It can be selected by name or path, loaded safely, and
expressed through a compact set of operations covering node identity, type,
interfaces, addressing, gateways, static routes, custom startup behaviour, and
OSPF, RIP, or BGP configuration. Structural validation protects the framework
from malformed input, while semantic validation examines the relationships
among those fields and distinguishes blocking errors from non-blocking
warnings. The result is a validated Python representation that can be handed
to the generation stage with a substantially lower risk of deterministic
configuration errors. The next chapter describes how this validated source is
translated into standard Kathar\'a artefacts and how the resulting laboratory
is managed during its lifecycle.
